% ==============================================================================
% reports/chapters/04_phuong_phap_de_xuat.tex
% CHƯƠNG 4: PHƯƠNG PHÁP VÀ KIẾN TRÚC HỆ THỐNG
% ==============================================================================

\chapter{Phương Pháp và Kiến Trúc Hệ Thống}
\label{chap:phuong_phap}

\section{Kiến trúc tổng thể của hệ thống thực nghiệm}

Nhằm tổ chức quy trình thực nghiệm và giảm một số nguy cơ rò rỉ thông tin giữa các giai đoạn, hệ thống được thiết kế theo đường ống gồm 5 giai đoạn; các bước này không tự thân chứng minh tính độc lập hoàn toàn giữa các tập:
\begin{enumerate}
    \item \textbf{Giai đoạn 1: Chuẩn bị dữ liệu và phân tách}: Áp dụng tái dựng timestamp và Purge/Embargo để chia dữ liệu thành Train, Validation và Test theo giao thức thời gian đã xác định.
    \item \textbf{Giai đoạn 2: Tiền xử lý chỉ học từ tập Train}: Loại bỏ các đặc trưng hằng số và khớp bộ điền giá trị thiếu theo trung vị duy nhất trên $\Dtrain$.
    \item \textbf{Giai đoạn 3: Tinh chỉnh siêu tham số}: Tìm kiếm cấu hình tối ưu của 3 mô hình cây trên không gian $\Dtrain \times \Dval$.
    \item \textbf{Giai đoạn 4: Đóng băng cấu hình}: Khóa tham số và cấu hình tốt nhất vào tệp siêu dữ liệu \texttt{frozen.json}.
    \item \textbf{Giai đoạn 5: Mở kiểm thử một lần duy nhất}: Đánh giá hiệu năng tổng quát hóa trên $\Dtest$ mà không tinh chỉnh lại bất kỳ tham số nào.
\end{enumerate}

\section{Chính sách quản lý đặc trưng}

Để ngăn chặn hiện tượng học lối tắt \cite{geirhos2020shortcut}, các trường thông tin định danh và siêu dữ liệu không phản ánh hành vi mạng thuần túy đều bị loại bỏ khỏi không gian đặc trưng:
\begin{itemize}
    \item \textbf{Các trường cấm sử dụng}: \texttt{Flow ID}, \texttt{Source IP}, \texttt{Destination IP}, \texttt{Source Port}, \texttt{Timestamp}, \texttt{Timestamp\_fixed}, \texttt{Protocol}, \texttt{TimeBin10}, \texttt{\_row\_id}, \texttt{index}. Việc giữ lại địa chỉ IP hay nhãn thời gian sẽ khiến mô hình ghi nhớ máy chủ tấn công thay vì học mẫu lưu lượng.
\end{itemize}

Chi tiết chính sách quản lý đặc trưng được trình bày trong Bảng \ref{tab:feature_schema}.

% ==============================================================================
% reports/tables/manual/tab_feature_schema.tex
% ==============================================================================
\begin{table}[htbp]
\centering
\small
\caption{Chính sách kiểm soát không gian đặc trưng và danh mục đặc trưng hằng số bị loại bỏ.}
\label{tab:feature_schema}
\begin{tabular}{lp{10cm}}
\toprule
\textbf{Nhóm đặc trưng} & \textbf{Chi tiết các trường thông tin và quy tắc xử lý} \\
\midrule
\textbf{Trường cấm tuyệt đối} & \texttt{Flow ID}, \texttt{Source IP}, \texttt{Destination IP}, \texttt{Source Port}, \texttt{Timestamp}, \texttt{Timestamp\_fixed}, \texttt{Protocol}, \texttt{TimeBin10}, \texttt{\_row\_id}, \texttt{index}. \\
\textbf{Lý do cấm} & Ngăn chặn mô hình học thuộc lòng địa chỉ máy chủ tấn công hoặc ghi nhớ thời điểm diễn ra chiến dịch. \\
\midrule
\textbf{10 Đặc trưng hằng số} & \texttt{Bwd PSH Flags}, \texttt{Fwd URG Flags}, \texttt{Bwd URG Flags}, \texttt{CWE Flag Count}, \texttt{Fwd Avg Bytes/Bulk}, \texttt{Fwd Avg Packets/Bulk}, \texttt{Fwd Avg Bulk Rate}, \texttt{Bwd Avg Bytes/Bulk}, \texttt{Bwd Avg Packets/Bulk}, \texttt{Bwd Avg Bulk Rate}. \\
\textbf{Cơ chế phát hiện} & Toàn bộ giá trị có phương sai bằng 0 trên tập Train ($\Dtrain$) $\rightarrow$ Loại bỏ đồng bộ trên Val và Test. \\
\midrule
\textbf{Không gian cuối cùng} & 67 đặc trưng (kịch bản \texttt{with\_port}) và 66 đặc trưng (kịch bản \texttt{without\_port}). \\
\bottomrule
\end{tabular}
\end{table}


\subsection{Thực nghiệm bóc tách cổng dịch vụ}
Cổng dịch vụ đích có giá trị bằng 21 cho giao thức FTP và 22 cho giao thức SSH. Để trả lời câu hỏi nghiên cứu RQ3 xem mô hình có bị phụ thuộc thái quá vào số cổng hay không, nhóm thiết kế hai nhánh không gian đặc trưng song song:
\begin{enumerate}
    \item \textbf{Nhánh có cổng} (\texttt{with\_port}, \FeatWithPortCount\ đặc trưng): Bao gồm trường \texttt{Destination Port} kết hợp với toàn bộ các đặc trưng thống kê hành vi luồng mạng.
    \item \textbf{Nhánh không có cổng} (\texttt{without\_port}, \FeatWithoutPortCount\ đặc trưng): Loại bỏ hoàn toàn trường \texttt{Destination Port}, buộc mô hình chỉ dựa vào các đặc trưng thống kê gói tin như kích thước, khoảng thời gian IAT, cờ TCP và độ dài tiêu đề gói tin.
\end{enumerate}

Danh mục định nghĩa chi tiết của toàn bộ \FeatWithPortCount\ đặc trưng mạng được cung cấp tại Phụ lục \ref{app:feature_list}.

\section{Thiết kế chiến lược phân tách dữ liệu}

\subsection{Cơ chế Purge và Embargo giảm nguy cơ vượt ranh giới thời gian}
Trong phân tích chuỗi thời gian mạng, một flow có thời điểm bắt đầu $s$ và thời lượng $d$ (tính bằng micro-giây) sẽ kéo dài qua một khoảng thời gian. Để tính thời điểm kết thúc theo cách bảo thủ, thời lượng được đổi sang giây và cộng thêm sai số làm tròn $q = 60\text{s}$:
\begin{equation}
    e = s + \frac{d}{10^6}\,\text{s} + q
\end{equation}

Tại mỗi mốc cắt $T_{\text{cutoff}}$, các flow có khoảng tồn tại vượt qua ranh giới thời gian sẽ bị loại bỏ:
\begin{enumerate}
    \item \textbf{Purge}: Loại bỏ flow bắt đầu trước $T_{\text{cutoff}}$ nhưng có thời điểm kết thúc bảo thủ $e \ge T_{\text{cutoff}}$.
    \item \textbf{Embargo}: Loại bỏ flow bắt đầu trong khoảng $[T_{\text{cutoff}}, T_{\text{cutoff}} + g)$, với khoảng đệm $g = 120\text{s}$.
\end{enumerate}

Các tham số $q = 60\text{s}$ và $g = 120\text{s}$ là lựa chọn thiết kế nhằm giảm nguy cơ flow vượt qua ranh giới thời gian hoặc có phản hồi trễ sát mốc cắt. Purge và Embargo không xác minh độ chính xác của đồng hồ nguồn, cũng không chứng minh tính độc lập hoàn toàn giữa các tập dữ liệu.

Sau khi khảo sát diễn biến nhãn và thời gian, nghiên cứu chốt hai mốc cắt:
\begin{itemize}
    \item \textbf{Mốc A}: \texttt{2017-07-04 10:00:00} (Phân cách Train và Validation)
    \item \textbf{Mốc B}: \texttt{2017-07-04 14:30:00} (Phân cách Validation và Test)
\end{itemize}

Thuật toán Purge và Embargo đã loại bỏ tổng cộng \textbf{\DataPurgedRows\ dòng} tại các vùng ranh giới (xem Bảng \ref{tab:dataset_splits}), qua đó giảm nguy cơ flow vượt qua các mốc phân tách thời gian.

\subsection{Đối chiếu hai giao thức phân tách dữ liệu}
Phân vùng Random Split được tạo bằng cách chia ngẫu nhiên phân tầng các flow, với Train ($20.6\%$), Validation ($46.5\%$) và Test ($32.9\%$). Kết quả của Random Split được đối chiếu với Time-based Split như hai giao thức đánh giá khác nhau, không phải phép thử cô lập riêng tác động của thuật toán chia. Trong Time-based Split, Train không có SSH và Test chỉ có SSH; trong Random Split, Train có 1,215 mẫu SSH và Test có 1,938 mẫu SSH cùng 2,610 mẫu FTP. Do thành phần subtype và phân bố lớp trong Test cũng khác nhau, chênh lệch điểm số phản ánh đồng thời các khác biệt này và chiến lược phân tách. Các flow phụ thuộc hoặc cùng chiến dịch xuất hiện ở nhiều tập là một cơ chế có thể góp phần làm ước lượng lạc quan, nhưng kết quả hiện có không trực tiếp chứng minh rò rỉ nhãn hay đặc trưng.

\section{Quy trình tiền xử lý độc lập trên tập train}

Một nguyên tắc quan trọng trong thiết kế đường ống MLOps là bộ tiền xử lý không được học tham số từ dữ liệu tương lai. Toàn bộ các bước học tham số tiền xử lý được thực hiện trên tập $\Dtrain$ theo quy trình chuẩn hóa:
\begin{enumerate}
    \item Loại bỏ đặc trưng hằng số: Thuật toán kiểm tra phương sai của từng đặc trưng trên $\Dtrain$. Có \FeatConstantDropped\ đặc trưng có giá trị không đổi (phương sai bằng 0) trên $\Dtrain$ (bao gồm các cờ hiếm như \texttt{Bwd PSH Flags}, \texttt{Fwd URG Flags}, \texttt{CWE Flag Count} và các trường bulk). Danh sách \FeatConstantDropped\ cột này được lưu lại và áp dụng loại bỏ đồng bộ trên $\Dval$ và $\Dtest$.
    \item Điền giá trị thiếu: Khởi tạo \texttt{SimpleImputer} với chiến lược trung vị (\texttt{strategy='median'}). Giá trị trung vị của từng đặc trưng được tính toán duy nhất trên $\Dtrain$, sau đó được dùng để điền các giá trị NaN trên cả ba tập $\Dtrain$, $\Dval$ và $\Dtest$.
\end{enumerate}

\section{Nguyên tắc đóng băng mô hình}

Để hạn chế nguy cơ tinh chỉnh siêu tham số lặp đi lặp lại trên tập Test (rò rỉ thông tin tập kiểm thử qua quá trình thử nghiệm), nhóm áp dụng quy tắc mở kiểm thử một lần duy nhất:
\begin{itemize}
    \item Toàn bộ quá trình thử nghiệm, tìm kiếm siêu tham số và lựa chọn mô hình đại diện được thực hiện hoàn toàn trên tập $\Dtrain$ và $\Dval$.
    \item Cấu hình tốt nhất sau đó được xuất ra tệp cấu hình \texttt{frozen.json} và tính mã băm SHA-256 xác thực.
    \item Quy trình đánh giá trên tập kiểm thử kiểm tra cờ xác thực \texttt{TEST\_OPENED.json}. Tập kiểm thử chỉ được mở đúng một lần duy nhất để ghi nhận chỉ số cuối cùng mà không tinh chỉnh lại mô hình hay tham số.
\end{itemize}
